\chapterhead{2}{Drying Up}{Water leaving as a gas you cannot see}\label{ch2}

Wet washing dries on a line. A wet towel dries on a hook. Nobody sees the water go. It leaves as \textbf{water vapour}, which is water as an invisible gas mixed into the air. Water turning into vapour is called \textbf{evaporation}.

The covered cup in Chapter~1 kept its water because the cling film kept the vapour in. The open cup let the vapour drift away.

\begin{tcolorbox}[colback=sun!14,colframe=sun!85!black,boxrule=1.2pt,arc=2mm,left=6pt,right=6pt,top=4pt,bottom=4pt,before skip=8pt,after skip=8pt]
\sans\bfseries\color{ink} Vapour is invisible. Steam and mist are not vapour.\par
\normalfont\sffamily\small The white cloud above a kettle, and a mist over a pond, are made of tiny drops of liquid water. Real vapour has no colour and cannot be seen. You will meet the tiny drops again in Chapter~3.
\end{tcolorbox}

\begin{tryit}{A drying race}
\textbf{You need:} two paper towel squares the same size, two plates, a spoon, a clock.
\begin{enumerate}[leftmargin=1.6em]
\item Put 4 spoonfuls of water on each square, so they are equally wet.
\item Lay one square \textbf{flat} on a plate. Fold the other into four and put it on a plate.
\item Put both plates side by side in the same spot. Check every 10 minutes. A square is dry when it feels dry and no dark patch is left.
\item Write the number of minutes each took. Do the race three times.
\end{enumerate}
\par\medskip
\renewcommand{\arraystretch}{1.6}
\begin{tabular}{|l|c|c|c|c|}
\hline
 & \textbf{Race 1} & \textbf{Race 2} & \textbf{Race 3} & \textbf{Average}\\ \hline
\textbf{Flat (min)} & & & & \\ \hline
\textbf{Folded (min)} & & & & \\ \hline
\end{tabular}
\end{tryit}

\begin{worked}{Practice numbers, not from a real test}
Flat: 40, 44, 42 minutes. Total 126, average $126\div3=42$, spread 4.\\
Folded: 130, 140, 135 minutes. Total 405, average $405\div3=135$, spread 10.\\
The folded square took $135\div42\approx3.2$, so about three times as long. The gap between the averages is much bigger than either spread.
\end{worked}

\begin{center}
\begin{tikzpicture}
  \draw[line width=1.2pt] (0,0) -- (0,3.6); \draw[line width=1.2pt] (0,0) -- (6.2,0);
  \foreach \y/\l in {0/0,0.9/40,1.8/80,2.7/120}{\draw (0,\y) -- (-0.1,\y) node[left,font=\sffamily\scriptsize] {\l};}
  \fill[sky!70,draw=ink] (1.0,0) rectangle (2.4,0.945); \node[font=\sffamily\scriptsize] at (1.7,1.25) {42};
  \fill[tangerine!80,draw=ink] (3.6,0) rectangle (5.0,3.04); \node[font=\sffamily\scriptsize] at (4.3,3.35) {135};
  \node[font=\sffamily\scriptsize] at (1.7,-0.35) {flat}; \node[font=\sffamily\scriptsize] at (4.3,-0.35) {folded};
  \node[font=\sffamily\scriptsize,rotate=90] at (-0.9,1.8) {minutes to dry};
\end{tikzpicture}
\end{center}

\begin{multix}[hmm]\multixbot[hmm]\tcblower
Water only evaporates when it is hot. On a cold day my towel will never dry. I am very sure about this. Should I be?
\end{multix}

\trust{We wet each square with the same amount, used the same paper, the same spot and the same clock, and repeated the race three times.}
{The three flat times were all near 42 and the three folded times all near 135, so the difference is far larger than the wobble between repeats.}
{Whether it holds for cloth, for other rooms, or for water that is not spread over paper. We changed only how the paper lay.}

\begin{keepbox}
We reduced drying to a number of minutes. \textbf{We kept} \wline[3.4cm] \textbf{and we set aside} \wline[3.4cm].
\end{keepbox}

\clearpage

\subhead{Warm and cool}
Is evaporation faster in the warm? Test it with the \emph{same} kind of towel square, the same amount of water and the same way of lying. Change only the place.

\begin{tryit}{Sun or shade?}
Put one flat wet square on a plate on a sunny window sill and one on a plate on a shady shelf. Check every 10 minutes. Do it three times.\par\medskip
\renewcommand{\arraystretch}{1.6}
\begin{tabular}{|l|c|c|c|c|}
\hline
 & \textbf{Try 1} & \textbf{Try 2} & \textbf{Try 3} & \textbf{Average}\\ \hline
\textbf{Sunny (min)} & & & & \\ \hline
\textbf{Shady (min)} & & & & \\ \hline
\end{tabular}
\end{tryit}

\begin{worked}{Practice numbers}
Sunny: 54, 60, 57 minutes. Total 171, average 57, spread 6.\\
Shady: 96, 102, 99 minutes. Total 297, average 99, spread 6.\\
The shady square took $99-57=42$ minutes longer.
\end{worked}

\subhead{What makes it quicker?}
\begin{center}
\sffamily\small\renewcommand{\arraystretch}{1.5}
\begin{tabular}{@{}p{3.4cm}p{6.6cm}@{}}
\toprule
\textbf{Change} & \textbf{Evaporation is usually quicker when}\\
\midrule
Warmth & the water and air are warmer\\
Spread & the water is spread thin over a bigger area\\
Moving air & a breeze carries the vapour away\\
Dry air & the air already holds little vapour\\
\bottomrule
\end{tabular}
\end{center}

Notice that one fair test changes one of these at a time. Folding the square changed the spread. Sun against shade changed the warmth.

\trust{We repeated each placement three times and compared the averages.}
{The sunny times were all below the shady ones, by far more than the spread of six minutes.}
{That warmth was the only cause. A sunny sill may also be breezier or drier than a shelf.}

\begin{keepbox}
We said the sun makes things dry faster. \textbf{We kept} \wline[3.4cm] \textbf{and we set aside} \wline[3.4cm].
\end{keepbox}

\subhead{Add a breeze}
\begin{tryit}{Still air or fanned?}
Dry two flat wet squares the same way as before, one in still air and one with a book or card waved gently over it from a short way off. Keep the plates and place the same. Do it three times.\par\medskip
\renewcommand{\arraystretch}{1.6}
\begin{tabular}{|l|c|c|c|c|}
\hline
 & \textbf{Try 1} & \textbf{Try 2} & \textbf{Try 3} & \textbf{Average}\\ \hline
\textbf{Still air (min)} & & & & \\ \hline
\textbf{Fanned (min)} & & & & \\ \hline
\end{tabular}
\end{tryit}

\begin{worked}{Practice numbers}
Still air: 42, 40, 44 minutes. Total 126, average 42, spread 4.\\
Fanned: 24, 28, 26 minutes. Total 78, average 26, spread 4.\\
The fanned square dried $42-26=16$ minutes sooner. A breeze carries vapour away, so more can leave.
\end{worked}

\begin{multix}[ok]\multixbot[ok]\tcblower
That means washing dries well on a windy day, even a cool one. Now my wet towel and I have something to talk about.
\end{multix}

\clearpage

\begin{puzzle}{Average the races}
Flat squares dried in 36, 40 and 38 minutes. Folded squares dried in 120, 126 and 132 minutes.\\
Flat average: \wline[1.6cm] \quad Folded average: \wline[1.6cm]
\end{puzzle}

\begin{puzzle}{Which dries first?}
Which of these dries first? Write 1 for first and 3 for last.\\
A towel on a line in a breeze \wline[1cm] \quad the same towel in a heap on the floor \wline[1cm] \quad the same towel folded in a closed drawer \wline[1cm]
\end{puzzle}

\begin{puzzle}{Spot the unfair test}
Sam dried a small square in the sun and a big square in the shade. The small one dried first. He says sun makes drying faster. List two things that changed. \wline[6cm]
\end{puzzle}

\begin{puzzle}{How many times as long?}
The folded square took 135 minutes and the flat one 42 minutes. About how many times as long did the folded square take? Round to a whole number. \wline[1.6cm]
\end{puzzle}

\begin{puzzle}{Vapour or drops?}
Write V if it is vapour and D if it is tiny drops.\\
(a) The cloud above a kettle \wline[1cm] \quad (b) The water in dry air that you cannot see \wline[1cm] \quad (c) Mist on a pond \wline[1cm]
\end{puzzle}

\clearpage

\begin{puzzle}{Wash day}
Washing dries on a cold, windy day more slowly than on a warm, windy one. It still dries. What does that tell you about Multix's claim? \wline[6cm]
\end{puzzle}

\begin{puzzle}{Move the water}
A towel holds 4 spoonfuls of water. After it is dry, where is the 4 spoonfuls? Fill in the circles.
\begin{center}
\begin{tikzpicture}
  \node[cir] (a) at (0,0) {4}; \node[cir,dashed] (b) at (5,0) {};
  \draw[moveA] (a) -- node[above,note] {evaporation} (b);
  \node[note] at (0,-0.8) {towel before}; \node[note] at (5,-0.8) {air gained};
\end{tikzpicture}
\end{center}
\end{puzzle}

\begin{puzzle}{Why folded is slower}
Say in one sentence why a folded towel dries more slowly than a flat one. \wline[6cm]
\end{puzzle}

\begin{puzzle}{Spread out}
Which dries faster: a cup of water, or the same water in a wide tray? \wline[3cm]\\
Why? \wline[5cm]
\end{puzzle}

\begin{puzzle}{Make a test}
Plan a fair test of one thing that might change how fast a towel dries. Change: \wline[3cm] Keep the same: \wline[4cm] Measure: \wline[3cm]
\end{puzzle}
\clearpage
